%======================================================================%
\section{Complete Standard Model Lagrangian}
\label{sec:sm-lagrangian}
%======================================================================%

This section records the \textbf{Standard Model Lagrangian certificate}:
every operator in the low-energy effective theory contacted at the
$\mathrm{SU}(3)_c\times\mathrm{SU}(2)_L\times U(1)_Y$ vacuum of
\eqref{eq:breaking-chain} is the $\Pi_{\sigma=-1}$-projected image of the
void--mirror pipeline
\begin{equation}\label{eq:sm-pipeline}
  \Void,\;\Mirror
  \;\longrightarrow\;
  \mathfrak{U}^\star=\Climb^\infty(\Void)
  \;\longrightarrow\;
  \hat Q=Q/\sigma
  \;\longrightarrow\;
  \mathcal{L}_{\mathrm{SM}},
\end{equation}
with no independent SM input beyond the two primitive axioms and the derived
assumption ledger \textbf{A1--A8}
(Definition~\ref{def:assumption-ledger}). The EIII coset
$\mathcal{M}=\E{6}/H$ supplies the kinetic metric $K$ and coercive potential
$V$ (Definitions~\ref{def:potential},~\eqref{eq:action}); sequential
$\iota$-ordered breaking (Theorem~\ref{thm:breaking-complete}) fixes gauge
content, Higgs multiplets, and Yukawa portals; observer collapse
(Definition~\ref{def:physical}) restricts all observable fields to
$H^\bullet(\hat Q,V_{16})^{\sigma=-1}$ on the mirror quotient
$\hat Q$ (Definition~\ref{def:Qhat}).

%----------------------------------------------------------------------%
\subsection{E6 coset, breaking chain, and SM gauge content}
\label{sec:sm-lagrangian:breaking}
%----------------------------------------------------------------------%

At the stabilized exceptional rung $\Mirror^8(\Void)=\E{6}$
(Theorem~\ref{thm:e6}), the universe datum carries the compact coset
$\mathcal{M}=\E{6}/(\Spin(10)\times U(1))$ with matter directions
$\mathfrak{m}=\mathbf{16}_{-3}\oplus\overline{\mathbf{16}}_{+3}$
(Theorem~\ref{thm:universe}). Observer collapse at the heterotic rung
$\Mirror^7(\Void)=\E{8}\times\E{8}$ selects the anti-invariant half
(Definition~\ref{def:physical}); the Stone lift and quotient
$\hat Q=Q/\sigma$ transport the coset dynamics to the continuum
configuration space used below (Theorem~\ref{thm:duality},
Proposition~\ref{prop:stone-lift}).

\begin{proposition}[SM gauge algebra from staged breaking]
\label{prop:sm-gauge-content}
Assume Theorem~\ref{thm:breaking-complete}. After the four capacity-ordered
stages of \eqref{eq:breaking-chain}, the residual gauge group on the
soldered $d=4$ block $H_2(\Quat)$ (Theorem~\ref{thm:spacetime:soldering-derived})
is
\begin{equation}\label{eq:sm-gauge-group}
  G_{\mathrm{SM}}
  \;=\;
  \mathrm{SU}(3)_c\times\mathrm{SU}(2)_L\times U(1)_Y,
\end{equation}
with Lie algebra $\mathfrak{g}_{\mathrm{SM}}=\mathfrak{su}_3\oplus\mathfrak{su}_2
\oplus\mathfrak{u}_1$. The eight gluons, three weak bosons
$W^\pm,Z^0$, and photon arise as the $\Pi_{\sigma=-1}$-projected massless
connections along the unbroken directions of \eqref{eq:breaking-chain}; the
$W$ and $Z$ masses are generated by electroweak flow along $\Phi_{\mathbf{10}}$
(Theorem~\ref{thm:breaking-ew}).
\end{proposition}

\begin{proof}
Theorem~\ref{thm:breaking-complete} realizes
\eqref{eq:breaking-chain} by gradient flow of $V_{\mathrm{eff}}$ on
$\iota$-restricted coset directions. Stage~(3) leaves
$\mathrm{SU}(3)_c\times\mathrm{SU}(2)_L\times U(1)_Y$
(Theorem~\ref{thm:breaking-su5sm}); stage~(4) breaks
$\mathrm{SU}(2)_L\times U(1)_Y\to U(1)_{\mathrm{em}}$ while preserving
$\mathrm{SU}(3)_c$ (Theorem~\ref{thm:breaking-ew}). Standard $\mathrm{SU}(5)$
branching under $\Pi_{\sigma=-1}$ (Lemma~\ref{lem:breaking-126-sigma})
assigns the adjoint $\mathbf{24}$ of $\mathrm{SU}(5)$ to
$(8,1)_0\oplus(1,3)_0\oplus(1,1)_0$ under
$\mathrm{SU}(3)_c\times\mathrm{SU}(2)_L\times U(1)_Y$, yielding the SM gauge
multiplet. Massless gauge bosons correspond to unbroken generators; massive
$W/Z$ arise from $\langle\Phi_{\mathbf{10}}\rangle$ in the $\mathbf{10}_H$
channel (Proposition~\ref{prop:breaking-higgs-labels}).
\end{proof}

%----------------------------------------------------------------------%
\subsection{Field content and covariant derivatives}
\label{sec:sm-lagrangian:fields}
%----------------------------------------------------------------------%

Physical fermions are cohomology classes in
$H^1(\hat Q,V_{16})^{\sigma=-1}$ (Definition~\ref{def:physical}), with
$\dim H^1(\hat Q,V_{16})^{\sigma=-1}=3$ observer slots
(Theorem~\ref{thm:chirality-derived}). Each slot carries one chiral
$\Spin(10)$ spinor $\mathbf{16}$, which branches to one SM generation after
\eqref{eq:breaking-chain}.

\begin{definition}[SM field content on $\hat Q$]\label{def:sm-fields}
Fix the $\mathrm{SU}(3)_c\times\mathrm{SU}(2)_L\times U(1)_Y$ vacuum of
Theorem~\ref{thm:breaking-complete}. The \emph{observable SM field content}
on $\hat Q$ is:
\begin{enumerate}[label=\textbf{(\arabic*)},leftmargin=*]
  \item \textbf{Gauge fields:} $A_\mu^a$ for $a=1,\ldots,8$ ($G_\mu^a$, gluons),
    $a=9,10,11$ ($W_\mu^a$), and $B_\mu$ ($U(1)_Y$), together spanning the
    adjoint of $G_{\mathrm{SM}}$ in \eqref{eq:sm-gauge-group}.
  \item \textbf{Three chiral generations} ($\iota=1,2,3$): for each observer
    index, left-handed quark doublets $Q_L^\iota=(u_L^\iota,d_L^\iota)$,
    right-handed up-type $u_R^\iota$, right-handed down-type $d_R^\iota$,
    left-handed lepton doublets $L_L^\iota=(\nu_L^\iota,e_L^\iota)$, and
    right-handed charged leptons $e_R^\iota$, realized as components of
    $|\omega_\iota\rangle\in H^1(\hat Q,V_{16})^{\sigma=-1,\,\iota}$
    under standard $\mathrm{SU}(5)\subset\Spin(10)$ branching.
  \item \textbf{Higgs doublet:} one complex $\mathrm{SU}(2)_L$ doublet
    $H=(H^+,H^0)$ from the $\Pi_{\sigma=-1}$-projected $\mathbf{10}_H$
    fluctuation $\Phi_{\mathbf{10}}$ with electroweak VEV
    $\langle H^0\rangle=v/\sqrt{2}$ (Theorem~\ref{thm:breaking-ew},
    Theorem~\ref{thm:yukawa-complete}).
  \item \textbf{No mirror partners:} $\overline{\mathbf{16}}$ modes in
    $H^1(\hat Q,V_{16})^{\sigma=+1}$ are quotient-non-observable
    (Theorem~\ref{thm:duality}); no fourth generation
    (Proposition~\ref{prop:no-fourth}).
\end{enumerate}
The SM covariant derivative on a field $\psi$ of hypercharge $Y$ is
\begin{equation}\label{eq:sm-covariant}
  D_\mu\psi
  \;=\;
  \partial_\mu\psi
  \;-\;
  i\,g_s\,G_\mu^a\,T^a\psi
  \;-\;
  i\,g\,W_\mu^a\,\tau^a\psi
  \;-\;
  i\,g'\,B_\mu\,Y\psi,
\end{equation}
with $T^a$ the $\mathrm{SU}(3)_c$ generators, $\tau^a$ the $\mathrm{SU}(2)_L$
generators, and $(g_s,g,g')$ the low-energy gauge couplings descended from the
$\Spin(10)$ normalization of Proposition~\ref{prop:weinberg-gut} and the
one-loop flow of Theorem~\ref{thm:beta} ($b_0=12$).
\end{definition}

\begin{remark}[Tier assignment for field content]
Items~\textbf{(1)--(3)} are Tier~B: gauge quantum numbers follow from
Theorem~\ref{thm:breaking-complete} (ledger entries \textbf{A3}, \textbf{A6});
three generations are Tier~A
(Theorem~\ref{thm:chirality-derived}). Absolute coupling values are Tier~C
contacts (Section~\ref{sec:scales}, step~\textbf{D3}).
\end{remark}

%----------------------------------------------------------------------%
\subsection{Gauge kinetic terms}
\label{sec:sm-lagrangian:gauge}
%----------------------------------------------------------------------%

The EIII $\sigma$-model \eqref{eq:action} is defined on the coset
$\mathcal{M}$ with $H$-invariant metric $K$ (Proposition~\ref{prop:coset-metric}).
Gauge connections on $\hat Q$ extend the coset tangent transport to the
$\Spin(10)\times U(1)$ bundle (Section~\ref{sec:physics:gauge}); after
\eqref{eq:breaking-chain}, the observable gauge kinetic terms are the
$\Pi_{\sigma=-1}$-projected Yang--Mills densities.

\begin{definition}[Gauge kinetic sector]\label{def:sm-gauge}
Let $F_{\mu\nu}^a$ denote the field strengths of $G_{\mathrm{SM}}$ in
\eqref{eq:sm-gauge-group}. The \emph{gauge kinetic sector} of the SM
Lagrangian is
\begin{equation}\label{eq:sm-gauge}
  \mathcal{L}_{\mathrm{gauge}}
  \;=\;
  -\frac{1}{4}\,G_{\mu\nu}^a G^{a,\mu\nu}
  \;-\;
  \frac{1}{4}\,W_{\mu\nu}^a W^{a,\mu\nu}
  \;-\;
  \frac{1}{4}\,B_{\mu\nu} B^{\mu\nu},
\end{equation}
with canonical normalization fixed by the $\Spin(10)$ embedding of
Proposition~\ref{prop:weinberg-gut} and the observer-local FRG
\eqref{eq:frg-system} at scale $\mu\equiv v$ for the electroweak contacts.
The low-energy values $(g_s,g,g')$ are RG images of the GUT normalization
$\sin^2\theta_W(M_{\mathrm{GUT}})=3/8$ (Proposition~\ref{prop:weinberg-gut}).
\end{definition}

\begin{lemma}[Gauge kinetic terms from $\Pi_{\sigma=-1}$]
\label{lem:sm-gauge-derived}
The Yang--Mills densities \eqref{eq:sm-gauge} are the $\Pi_{\sigma=-1}$-projected
quadratic fluctuation of the $H$-gauge connection on $\hat Q$ about the
$\mathrm{SU}(3)_c\times\mathrm{SU}(2)_L\times U(1)_Y$ vacuum of
Theorem~\ref{thm:breaking-complete}. Unbroken $\mathrm{SU}(3)_c$ directions
remain massless after stage~(4); $\mathrm{SU}(2)_L$ and mixed
$U(1)_Y$ directions acquire mass from $\langle\Phi_{\mathbf{10}}\rangle$.
\end{lemma}

\begin{proof}
Expand the gauge connection $A=A_{\mathrm{res}}+a$ about the residual vacuum.
$H$-invariance of \eqref{eq:action} and $\sigma$-equivariance
(Theorem~\ref{thm:sigma-measure}, Step~1) restrict quadratic gauge
fluctuations to the $\sigma=-1$ sector (Definition~\ref{def:physical}).
Theorem~\ref{thm:breaking-complete} identifies which adjoint directions are
unbroken at each stage; after stage~(4), only the eight gluons and photon are
exactly massless at tree level, with $W/Z$ masses from the Higgs mechanism
(Theorem~\ref{thm:breaking-ew}). The normalization follows from the
$\mathrm{SU}(5)\subset\Spin(10)$ embedding used in
Proposition~\ref{prop:weinberg-gut}.
\end{proof}

%----------------------------------------------------------------------%
\subsection{Fermion kinetic terms}
\label{sec:sm-lagrangian:fermion}
%----------------------------------------------------------------------%

\begin{definition}[Fermion kinetic sector]\label{def:sm-fermion}
For each generation $\iota\in\{1,2,3\}$, let $\psi_{\iota}$ denote the chiral
spinor components of Definition~\ref{def:sm-fields}. The \emph{fermion kinetic
sector} is
\begin{equation}\label{eq:sm-fermion}
  \mathcal{L}_{\mathrm{fermion}}
  \;=\;
  \sum_{\iota=1}^{3}\;
  \Bigl[
    i\,\bar Q_L^\iota\,\slashed{D}\,Q_L^\iota
    + i\,\bar u_R^\iota\,\slashed{D}\,u_R^\iota
    + i\,\bar d_R^\iota\,\slashed{D}\,d_R^\iota
    + i\,\bar L_L^\iota\,\slashed{D}\,L_L^\iota
    + i\,\bar e_R^\iota\,\slashed{D}\,e_R^\iota
  \Bigr],
\end{equation}
with $\slashed{D}=\gamma^\mu D_\mu$ and $D_\mu$ as in \eqref{eq:sm-covariant}.
The $\gamma$-matrices and spinor structure are imported by OS reconstruction on
$d\mu_{\Gamma_{-1}}$ (Proposition~\ref{prop:os-observer}, Theorem~\ref{thm:qm-derived})
together with soldering on $H_2(\Quat)$
(Theorem~\ref{thm:spacetime:soldering-derived}).
\end{definition}

\begin{lemma}[Fermion kinetic terms from $H^1(\hat Q,V_{16})^{\sigma=-1}$]
\label{lem:sm-fermion-derived}
Each chiral spinor in \eqref{eq:sm-fermion} is a $\Pi_{\sigma=-1}$-projected
Dolbeault--Serre representative of $|\omega_\iota\rangle$ on
$(\hat Q,V_{16})$ (Definition~\ref{def:observer}). The kinetic operator is the
observable restriction of the $H$-covariant Dirac operator coupled to the SM
connection \eqref{eq:sm-covariant}.
\end{lemma}

\begin{proof}
Definition~\ref{def:physical} declares physical fermions to be classes in
$H^1(\hat Q,V_{16})^{\sigma=-1}$; Theorem~\ref{thm:chirality-derived} gives
three such classes, one per observer slot. The Stone lift identifies cellular
cocycles on $\mathcal{T}_7$ with continuum $1$-forms on $\hat Q$
(Proposition~\ref{prop:stone-lift}). OS reconstruction
(Theorem~\ref{thm:qm-derived}) supplies the spinor bundle and Dirac operator on
the soldered $d=4$ manifold; coupling to $A_\mu^a$ yields \eqref{eq:sm-covariant}.
Mirror partners in $H^1(\hat Q,V_{16})^{\sigma=+1}$ do not appear
(Theorem~\ref{thm:duality}).
\end{proof}

%----------------------------------------------------------------------%
\subsection{Yukawa sector}
\label{sec:sm-lagrangian:yukawa}
%----------------------------------------------------------------------%

\begin{definition}[Yukawa sector]\label{def:sm-yukawa}
Let $H$ be the Higgs doublet of Definition~\ref{def:sm-fields} and let
$y_{ij}^{u}$, $y_{ij}^{d}$, $y_{ij}^{e}$ denote the $3\times3$ Yukawa matrices
generated by cubic overlap integrals \eqref{eq:breaking-yukawa}. The
\emph{Yukawa sector} is
\begin{equation}\label{eq:sm-yukawa}
  \mathcal{L}_{\mathrm{Yukawa}}
  \;=\;
  -\sum_{i,j=1}^{3}
  \Bigl[
    \bar Q_{L,i}\,y_{ij}^{d}\,H\,d_{R,j}
    + \bar Q_{L,i}\,y_{ij}^{u}\,\widetilde{H}\,u_{R,j}
    + \bar L_{L,i}\,y_{ij}^{e}\,H\,e_{R,j}
  \Bigr]
  + \mathrm{h.c.},
\end{equation}
with $\widetilde{H}=i\sigma_2 H^*$ the $\mathrm{SU}(2)_L$-conjugate doublet.
Up-type couplings use $\Phi_{\mathbf{126}}\in\mathbf{126}_H$ portal;
down-type and lepton couplings use $\Phi_{\mathbf{10}}\in\mathbf{10}_H$
(Proposition~\ref{prop:breaking-higgs-labels}, Theorem~\ref{thm:yukawa-complete}).
\end{definition}

\begin{lemma}[Yukawa matrices from overlap integrals]
\label{lem:sm-yukawa-derived}
The entries $y_{ij}$ in \eqref{eq:sm-yukawa} satisfy
\begin{equation}\label{eq:sm-yukawa-overlap}
  y_{ij}
  \;\propto\;
  \int_{\hat Q}\omega_i\wedge\omega_j\wedge\Phi_{\mathrm{Higgs}},
\end{equation}
with $\omega_i$ the representative of $|\omega_i\rangle$,
$\Phi_{\mathrm{Higgs}}\in\{\Phi_{\mathbf{10}},\Phi_{\mathbf{126}}\}$, and
leading-order factorization
$\mathcal{Y}_{ij}=y_0\sqrt{w_iw_j}\mathcal{P}_{ij}$
(Proposition~\ref{prop:yukawa-factor}, Theorem~\ref{thm:Hij-derived}).
\end{lemma}

\begin{proof}
Theorem~\ref{thm:breaking-complete} defines \eqref{eq:breaking-yukawa} on
$\hat Q$. Lemma~\ref{lem:yukawa-cellular} reduces the cubic overlap to
$\mathcal{T}_7$ hub pairings; Theorem~\ref{thm:yukawa-complete} closes the
full $3\times3$ structure with diagonal spectrum
Theorem~\ref{thm:yukawa-mass} and off-diagonal mixing from
Theorem~\ref{thm:Hij-derived}. Channel assignment
($u$-type vs.\ $d$-type/lepton) follows from the $\mathbf{126}_H$ vs.\
$\mathbf{10}_H$ labels in \eqref{eq:breaking-higgs-map}.
\end{proof}

%----------------------------------------------------------------------%
\subsection{Higgs potential}
\label{sec:sm-lagrangian:higgs}
%----------------------------------------------------------------------%

\begin{definition}[Higgs sector]\label{def:sm-higgs}
The \emph{Higgs kinetic and potential sector} is
\begin{equation}\label{eq:sm-higgs}
  \mathcal{L}_{\mathrm{Higgs}}
  \;=\;
  (D_\mu H)^\dagger(D^\mu H)
  \;-\;
  V_{\mathrm{Higgs}}(H),
  \qquad
  V_{\mathrm{Higgs}}(H)
  \;=\;
  -\mu_H^2\,H^\dagger H
  \;+\;
  \lambda_H\,(H^\dagger H)^2,
\end{equation}
with $\mu_H^2>0$ and $\lambda_H>0$ the $\Pi_{\sigma=-1}$-projected low-energy
parameters of the unstable $\Hess_1 V_{\mathrm{eff}}$ mode along
$\Phi_{\mathbf{10}}$ (Lemma~\ref{lem:breaking-cw},
Theorem~\ref{thm:breaking-ew}). The electroweak VEV is
$v^2=2\mu_H^2/\lambda_H$, identified with the $\iota=1$ capacity scale
(Theorem~\ref{thm:breaking-complete}, Section~\ref{sec:scales}).
\end{definition}

\begin{lemma}[Higgs potential from $V_{\mathrm{eff}}$]
\label{lem:sm-higgs-derived}
The quartic potential \eqref{eq:sm-higgs} is the $\Pi_{\sigma=-1}$-projected
Taylor expansion of $V_{\mathrm{eff}}$ about $\orderparam_\star$ along the
$\mathbf{10}_H$ fluctuation $\Phi_{\mathbf{10}}$, restricted to the physical
Higgs doublet component after electroweak symmetry breaking. The mass parameter
$\mu_H^2$ is the Coleman--Weinberg tachyonic coefficient on the $\iota=1$ cell
(Lemma~\ref{lem:breaking-cw}); $\lambda_H$ is the quartic coupling inherited
from the coercive $V$ of Definition~\ref{def:potential} at the EIII minimum.
\end{lemma}

\begin{proof}
Coercivity of $V$ (Theorem~\ref{thm:vacuum}) fixes a quartic tail on
$\mathcal{M}$; Coleman--Weinberg lifting (Lemma~\ref{lem:breaking-cw}) gaps
the $28$ coset directions and leaves one negative mode per $\iota$-cell. The
$\iota=1$ unstable direction maps to $\Phi_{\mathbf{10}}$
(Proposition~\ref{prop:breaking-higgs-labels}); gradient flow along $g_1$
realizes electroweak breaking (Theorem~\ref{thm:breaking-ew}). Expanding
$V_{\mathrm{eff}}$ in the $\mathrm{SU}(2)_L$ doublet component of
$\Phi_{\mathbf{10}}$ yields the standard Higgs potential with positive
$\lambda_H$ from the coercive background.
\end{proof}

%----------------------------------------------------------------------%
\subsection{Ghost sector and gauge fixing}
\label{sec:sm-lagrangian:ghost}
%----------------------------------------------------------------------%

BRST gauge fixing on the $H$-bundle over $\hat Q$ is required to define the
Euclidean measure \eqref{eq:gamma-factor} (Theorem~\ref{thm:sigma-measure}).
After \eqref{eq:breaking-chain}, Faddeev--Popov ghosts attach to the residual
$G_{\mathrm{SM}}$ connections.

\begin{definition}[Ghost sector]\label{def:sm-ghost}
Let $c^a$ and $\bar c^a$ denote the Faddeev--Popov ghost fields in the adjoint
of $G_{\mathrm{SM}}$. The \emph{ghost sector} is
\begin{equation}\label{eq:sm-ghost}
  \mathcal{L}_{\mathrm{ghost}}
  \;=\;
  \bar c^a\,\partial_\mu D^{ab} D^\mu c^b,
\end{equation}
with $D^{ab}=\delta^{ab}\partial_\mu-g_sf^{abc}G_\mu^c-g\epsilon^{abc}W_\mu^c
-g'd^{ab}B_\mu$ the gauge-covariant ghost kinetic operator in the SM vacuum.
Equivalently, the ghost fluctuation operator on $\hat Q$ is the adjoint
covariant Laplacian
\begin{equation}\label{eq:sm-ghost-operator}
  \Delta_{\mathrm{ghost}}
  \;=\;
  -\bigl(D_{\mathrm{adj}}\bigr)^2,
  \qquad
  D_{\mathrm{adj}}=d+[A,\cdot],
\end{equation}
acting on $\mathrm{adj}(G_{\mathrm{SM}})$; ghosts are massless and do not
contribute to the Sakharov $a_2$ sum (Theorem~\ref{thm:planck}).
\end{definition}

\begin{lemma}[Ghost sector from $\Pi_{\sigma=-1}$ gauge fixing]
\label{lem:sm-ghost-derived}
The ghost action \eqref{eq:sm-ghost} is the $\Pi_{\sigma=-1}$-projected
Faddeev--Popov determinant of local $G_{\mathrm{SM}}$ gauge fixing on
$\hat Q$. Ghost fields are Lorentz scalars in $\mathrm{adj}(G_{\mathrm{SM}})$;
their measure factor enters $d\mu_{\Gamma_{-1}}$ only through the observable
tensor factor of \eqref{eq:gamma-split}.
\end{lemma}

\begin{proof}
Theorem~\ref{thm:sigma-measure} factorizes the path integral after gauge
fixing on each $\sigma$-sector (Step~4). The $H$-covariant gauge fixing at the
$\Spin(10)\times U(1)$ level descends to $G_{\mathrm{SM}}$ along
\eqref{eq:breaking-chain}; $\Pi_{\sigma=-1}$ projects the adjoint ghost
multiplet to the observable gauge algebra of
Proposition~\ref{prop:sm-gauge-content}. Masslessness follows because ghosts
transform as adjoint scalars with no mass term in the BRST-exact sector; they
do not enter the $28$-scalar Sakharov induction (Theorem~\ref{thm:planck}).
\end{proof}

%----------------------------------------------------------------------%
\subsection{Strong CP: $\theta_{\mathrm{QCD}}=0$}
\label{sec:sm-lagrangian:theta}
%----------------------------------------------------------------------%

The QCD topological term is absent from the observable SM Lagrangian.

\begin{proposition}[Observable $\theta_{\mathrm{QCD}}=0$]
\label{prop:sm-theta-zero}
Assume Theorem~\ref{thm:strong-cp} and Definition~\ref{def:theta-qcd}. The
$\Pi_{\sigma=-1}$-projected SM effective action contains no
$\theta_{\mathrm{QCD}}\,\mathrm{Tr}(G\wedge G)$ insertion:
\begin{equation}\label{eq:sm-theta-zero}
  \theta_{\mathrm{obs}}
  \;=\;
  \Pi_{\sigma=-1}\,\theta_{\mathrm{QCD}}
  \;=\;
  0.
\end{equation}
The neutron EDM constraint is satisfied without a Peccei--Quinn axion field:
$\theta_{\mathrm{QCD}}$ is sequestered in the mirror-shadow sector
$d\mu_{\Gamma_{+1}}$ (Proposition~\ref{prop:theta-measure}).
\end{proposition}

\begin{proof}
Lemma~\ref{lem:theta-parity} shows $\theta_{\mathrm{QCD}}\,
\mathrm{Tr}(F\wedge F)$ is $\sigma$-odd; Proposition~\ref{prop:theta-measure}
gives $\theta_{\mathrm{obs}}=0$ after $\Pi_{\sigma=-1}$ projection.
Theorem~\ref{thm:strong-cp} records this as the Tier~B strong-$CP$ solution.
\end{proof}

%----------------------------------------------------------------------%
\subsection{Complete Lagrangian and derivation theorem}
\label{sec:sm-lagrangian:complete}
%----------------------------------------------------------------------%

\begin{definition}[Standard Model Lagrangian]\label{def:sm-lagrangian}
The \emph{Standard Model Lagrangian} on the soldered spacetime imported from
$H_2(\Quat)$ is the sum of the five sectors above:
\begin{equation}\label{eq:sm-lagrangian}
  \mathcal{L}_{\mathrm{SM}}
  \;=\;
  \mathcal{L}_{\mathrm{gauge}}
  \;+\;
  \mathcal{L}_{\mathrm{fermion}}
  \;+\;
  \mathcal{L}_{\mathrm{Yukawa}}
  \;+\;
  \mathcal{L}_{\mathrm{Higgs}}
  \;+\;
  \mathcal{L}_{\mathrm{ghost}},
\end{equation}
with explicit forms \eqref{eq:sm-gauge}, \eqref{eq:sm-fermion},
\eqref{eq:sm-yukawa}, \eqref{eq:sm-higgs}, and \eqref{eq:sm-ghost}. The
covariant derivative is \eqref{eq:sm-covariant}; the field content is
Definition~\ref{def:sm-fields}. The QCD $\theta$-term is omitted by
\eqref{eq:sm-theta-zero}. Anomaly cancellation is enforced by
Theorem~\ref{thm:anomaly}.
\end{definition}

\begin{table}[ht]
\centering
\small
\renewcommand{\arraystretch}{1.12}
\begin{tabular}{@{}llp{0.38\linewidth}c@{}}
\toprule
\textbf{SM operator} & \textbf{SJET source} & \textbf{Derivation route} & \textbf{Tier} \\
\midrule
$G_{\mu\nu}^a G^{a,\mu\nu}$ (gluons) &
  $K$ on $\mathfrak{m}$, $H$-bundle on $\hat Q$ &
  \eqref{eq:action}; Thm.~\ref{thm:breaking-complete} stage~(3)--(4);
  Lem.~\ref{lem:sm-gauge-derived} & B \\
$W_{\mu\nu}^a W^{a,\mu\nu}$, $B_{\mu\nu}B^{\mu\nu}$ &
  $\Phi_{\mathbf{10}}$ VEV, $\Spin(10)$ embedding &
  Thm.~\ref{thm:breaking-ew}; Prop.~\ref{prop:weinberg-gut} & B \\
$\bar\psi\slashed{D}\psi$ (3 gen.) &
  $H^1(\hat Q,V_{16})^{\sigma=-1}$ &
  Def.~\ref{def:physical}; Thm.~\ref{thm:chirality-derived};
  Lem.~\ref{lem:sm-fermion-derived} & A/B \\
$y_{ij}\bar\psi H\psi$ &
  Cubic overlap on $\hat Q$ &
  \eqref{eq:breaking-yukawa}; Thm.~\ref{thm:yukawa-complete};
  Lem.~\ref{lem:sm-yukawa-derived} & B \\
$(D_\mu H)^\dagger(D^\mu H)$ &
  $\Hess_1 V_{\mathrm{eff}}$ on $\iota=1$ cell &
  Prop.~\ref{prop:breaking-higgs-labels}; Lem.~\ref{lem:sm-higgs-derived} & B \\
$V_{\mathrm{Higgs}}(H)$ &
  Coercive $V$ + CW lift &
  Def.~\ref{def:potential}; Lem.~\ref{lem:breaking-cw};
  Thm.~\ref{thm:breaking-ew} & B \\
$\bar c D^2 c$ (ghosts) &
  FP gauge fixing on $\hat Q$ &
  Thm.~\ref{thm:sigma-measure}; Lem.~\ref{lem:sm-ghost-derived} & B \\
$\theta_{\mathrm{QCD}}\,\mathrm{Tr}(G\wedge G)$ &
  $\sigma$-odd topological density &
  Lem.~\ref{lem:theta-parity}; Thm.~\ref{thm:strong-cp};
  \eqref{eq:sm-theta-zero} & B \\
Anomaly cancellation &
  $\mathbf{27}$ arithmetic, $n_F=3$ &
  Thm.~\ref{thm:anomaly}; Thm.~\ref{thm:chirality-derived} & A \\
\bottomrule
\end{tabular}
\caption{Complete SM operator map: every term in
\eqref{eq:sm-lagrangian} is traced to the void--mirror pipeline through
$\Pi_{\sigma=-1}$ on $\hat Q$. Tier labels follow
Definition~\ref{def:proof-tiers}.}
\label{tab:sm-operator-map}
\end{table}

\begin{theorem}[SM Lagrangian derived from observer collapse]
\label{thm:sm-lagrangian-derived}
Assume the master pipeline \eqref{eq:master-pipeline} with
$\mathfrak{U}^\star=\Climb^\infty(\Void)$, observer collapse
$\Pi_{\sigma=-1}$ on $\hat Q$, and the staged vacuum of
Theorem~\ref{thm:breaking-complete}. Then the low-energy effective Lagrangian
on the soldered $d=4$ observer sector is \eqref{eq:sm-lagrangian}, with:
\begin{enumerate}[label=\textbf{(\roman*)},leftmargin=*]
  \item \textbf{Gauge content} $G_{\mathrm{SM}}$ of
    \eqref{eq:sm-gauge-group} (Proposition~\ref{prop:sm-gauge-content}).
  \item \textbf{Three chiral generations} from
    $\dim H^1(\hat Q,V_{16})^{\sigma=-1}=3$
    (Theorem~\ref{thm:chirality-derived}, Definition~\ref{def:sm-fields}).
  \item \textbf{One Higgs doublet} from $\Phi_{\mathbf{10}}\in\mathbf{10}_H$
    (Theorem~\ref{thm:breaking-ew}, Theorem~\ref{thm:yukawa-complete}).
  \item \textbf{Covariant derivatives} \eqref{eq:sm-covariant} from the
    residual $G_{\mathrm{SM}}$ connection on $\hat Q$.
  \item \textbf{Yukawa matrices} from cubic overlap integrals
    (Theorem~\ref{thm:yukawa-complete}, Lemma~\ref{lem:sm-yukawa-derived}).
  \item \textbf{Higgs potential} from $V_{\mathrm{eff}}$ along the
    $\iota=1$ unstable mode (Lemma~\ref{lem:sm-higgs-derived}).
  \item \textbf{Ghost sector} from $\Pi_{\sigma=-1}$ BRST gauge fixing
    (Lemma~\ref{lem:sm-ghost-derived}).
  \item \textbf{$\theta_{\mathrm{QCD}}=0$} in the observable sector
    (Proposition~\ref{prop:sm-theta-zero}, Theorem~\ref{thm:strong-cp}).
  \item \textbf{Anomaly cancellation} for the full three-family content
    (Theorem~\ref{thm:anomaly}).
\end{enumerate}
Every operator in Table~\ref{tab:sm-operator-map} appears in
\eqref{eq:sm-lagrangian} and is fixed by the pipeline
\eqref{eq:sm-pipeline}; no additional SM postulate is introduced.
\end{theorem}

\begin{proof}
\emph{Step 1 (coset and breaking).}
Void and mirror fix $\mathfrak{U}^\star$ with $\E{6}$ coset data
(Theorem~\ref{thm:universe}). The action \eqref{eq:action} with coercive
$V$ (Theorem~\ref{thm:vacuum}) and Coleman--Weinberg lift
(Lemma~\ref{lem:breaking-cw}) supplies $V_{\mathrm{eff}}$. Sequential
$\iota$-ordered gradient flow realizes \eqref{eq:breaking-chain}
(Theorem~\ref{thm:breaking-complete}), fixing $G_{\mathrm{SM}}$ and the Higgs
labels $\mathbf{45}_H$, $\mathbf{126}_H$, $\mathbf{10}_H$
(Proposition~\ref{prop:breaking-higgs-labels}). This proves
\textbf{(i)} and the Higgs doublet of \textbf{(iii)}.

\emph{Step 2 (observer collapse and fermions).}
Definition~\ref{def:physical} restricts observables to
$H^\bullet(\hat Q,V_{16})^{\sigma=-1}$. Theorem~\ref{thm:chirality-derived}
gives three chiral $\mathbf{16}$s; standard $\mathrm{SU}(5)$ branching yields
the SM multiplet assignment of Definition~\ref{def:sm-fields}, proving
\textbf{(ii)}. OS reconstruction (Theorem~\ref{thm:qm-derived}) with soldering
(Theorem~\ref{thm:spacetime:soldering-derived}) supplies spinors and
$\slashed{D}$, proving \textbf{(iv)} and Lemma~\ref{lem:sm-fermion-derived}.

\emph{Step 3 (gauge and ghosts).}
Theorem~\ref{thm:sigma-measure} factorizes the gauge-fixed measure on
$\hat Q$. Quadratic gauge fluctuations about the stage~(4) vacuum yield
\eqref{eq:sm-gauge} (Lemma~\ref{lem:sm-gauge-derived}); BRST ghosts follow
from the adjoint Laplacian \eqref{eq:sm-ghost-operator}
(Lemma~\ref{lem:sm-ghost-derived}), proving \textbf{(vii)}.

\emph{Step 4 (Yukawa and Higgs).}
Theorem~\ref{thm:breaking-complete} defines cubic overlaps
\eqref{eq:breaking-yukawa}; Theorem~\ref{thm:yukawa-complete} evaluates the
full $3\times3$ matrices, yielding \eqref{eq:sm-yukawa} and \textbf{(v)}.
The $\iota=1$ unstable mode of $V_{\mathrm{eff}}$ generates
$V_{\mathrm{Higgs}}$ (Lemma~\ref{lem:sm-higgs-derived}), proving
\textbf{(vi)}.

\emph{Step 5 ($\theta_{\mathrm{QCD}}$ and anomalies).}
Lemma~\ref{lem:theta-parity} and Theorem~\ref{thm:strong-cp} give
$\theta_{\mathrm{obs}}=0$ (Proposition~\ref{prop:sm-theta-zero}), proving
\textbf{(viii)}. Theorem~\ref{thm:anomaly} cancels gauge anomalies for
$n_F=3$, proving \textbf{(ix)}.

\emph{Step 6 (assembly).}
Definitions~\ref{def:sm-gauge}--\ref{def:sm-ghost} and
Definition~\ref{def:sm-lagrangian} assemble the five sectors into
\eqref{eq:sm-lagrangian}. Table~\ref{tab:sm-operator-map} records the
operator-level map with tier labels. No operator outside this list is required
for the observable SM EFT at scales below $M_{\mathrm{GUT}}$.
\end{proof}

\begin{corollary}[SM EFT closure]
\label{cor:sm-lagrangian-closure}
Theorem~\ref{thm:sm-lagrangian-derived} closes the \emph{field-theoretic}
description of the observable universe at electroweak scales: the SM
Lagrangian is a Tier~B consequence of void, mirror, observer collapse, and the
derived ledger \textbf{A1--A8}. Absolute scales $(v,\MPl)$ and sub-percent
flavor contacts are Tier~C programme steps~\textbf{D2--D3}
(Sections~\ref{sec:scales},~\ref{sec:yukawa}).
\end{corollary}

\begin{remark}[As above, so below on the Lagrangian]
The global EIII potential $V$ on $\mathfrak{U}^\star$ and the local SM
Lagrangian \eqref{eq:sm-lagrangian} on $\hat Q$ are linked by the
above--below correspondence (Definition~\ref{def:above-below},
Remark~\ref{rem:breaking-above-below}): capacity weights
$(1/27,10/27,16/27)$ govern both the profinite partition and the
$\iota$-ordered breaking thresholds that assemble \eqref{eq:sm-lagrangian}.
The Lagrangian is not postulated; it is the $\Pi_{\sigma=-1}$ effective
description of the mirror quotient at the heterotic rung.
\end{remark}